% §16.8 Notes and Lineage
%
% Monograph register: lineage of the concentration tools (Hoeffding, Azuma,
% McDiarmid, Devroye-Lugosi), the exact Bloom FPR (Bose et al. 2008;
% Christensen et al. 2010), adversarial filters (Naor-Yogev 2015;
% Clayton-Patton-Shrimpton 2019; Filic et al. 2022; Bender et al. 2018),
% adversarial ML (Goodfellow; Carlini-Wagner), hidden variables (Bell;
% Kochen-Specker; Fine), diagonal-POVM readout (Nielsen-Chuang; Bravyi et al.;
% Nation et al.; Smith et al.), cryptographic foundations (Goldreich).
% Replaces the former "Bibliographic Notes"; label kept. DEFINED:
% sec:bib-notes-ch16.

\section{Notes and Lineage}
\label{sec:bib-notes-ch16}

\paragraph{Lineage.}
None of the mathematics in this chapter is new. The concentration tools are
classical: \textcite{hoeffding-1963} for independent bounded summands,
\textcite{azuma-1967} for bounded martingale differences, and
\textcite{mcdiarmid-1989} for bounded-difference functions, with
\textcite{devroye-lugosi-2001} cataloguing their use on empirical
processes. That Bloom's false-positive formula is only a lower bound,
because the $k$ probes of one query are positively dependent, was proved
by \textcite{bose-et-al-2008}, and \textcite{christensen-roginsky-jimeno-2010}
gave the exact count. Adversarial use of filters has its own formal
literature: adversarial resilience \parencite{naor-yogev-2015},
correctness games under adaptive queries
\parencite{clayton-patton-shrimpton-2019}, their extension to privacy
\parencite{filic-et-al-2022}, and adaptive filters that repair false
positives \parencite{bender-et-al-2018}; adversarial examples for
classifiers trace to \textcite{goodfellow-2015} and
\textcite{carlini-wagner-2017}. The hidden-variable results are
\textcite{bell-1964}, \textcite{kochen-specker-1967}, and the equivalence
of \textcite{fine-1982}. That a diagonal POVM sees only the diagonal of
$\rho$ is the standard calculation \parencite{nielsen-chuang-2010}, and
the readout-error literature models computational-basis measurement with
exactly such assignment matrices
\parencite{bravyi-et-al-2021,nation-et-al-2021,smith-et-al-2021}.
\textcite{goldreich-1996} is the entry point to the cryptographic
primitives the chapter defers to.

\paragraph{What this chapter adds.}
An organization rather than a result: each failure mode is tied to the
assumption it violates (Axiom~2 for drift, Axiom~1 for shared randomness,
both for adaptive queries, and the prior existence of a joint distribution
for contextuality), and the replacement tool is named. Two statements of
earlier drafts are corrected here: the direction in which correlation
moves the serial-composition formula, with the Fr\'echet interval as the
only guarantee from marginals (\Cref{rem:composition-correlation-bracket}),
and the scope of the readout regime, which is fixed by the measurement
rather than by the state (\Cref{sec:quantum-contextuality}).

\paragraph{Left out.}
The security games and constructions of the adversarial-filter
literature, dependence calculi finer than the Fr\'echet interval,
readout-error mitigation, and any quantitative treatment of contextuality
beyond the Bell and Kochen-Specker obstructions.
